
\subsection{Grenzen der Schrödingergleichung}
\label{subsec:grenzen_schroedingergleichung}

Die Schrödingergleichung%
\index{Schrödingergleichung!Grenzen}
gehört zu den erfolgreichsten Gleichungen der modernen Physik. Mit ihr
lassen sich unter anderem Atome, Moleküle, chemische Bindungen und viele
Eigenschaften von Festkörpern beschreiben.

Trotz dieses Erfolges ist sie keine vollständige Theorie des Elektrons.
Sie gilt nur unter bestimmten Voraussetzungen. Vor allem ist sie eine
nichtrelativistische Gleichung%
\index{Quantenmechanik!nichtrelativistische}. Sie beschreibt Teilchen
zuverlässig, solange deren Geschwindigkeit deutlich kleiner als die
Lichtgeschwindigkeit ist.

Bei sehr hohen Energien treten jedoch relativistische Effekte auf.
Dann müssen Quantenmechanik und spezielle Relativitätstheorie%
\index{Relativitätstheorie!spezielle} miteinander verbunden werden.

\subsubsection*{Die nichtrelativistische Energie}

Die Schrödingergleichung beruht auf der klassischen Beziehung zwischen
kinetischer Energie und Impuls:

\begin{equation}
	E_{\mathrm{kin}}
	=
	\frac{p^2}{2m}.
	\label{eq:klassische_energie_schroedinger_grenzen}
\end{equation}

Dabei bezeichnet \(p\) den Impuls und \(m\) die Masse des Teilchens.

Diese Beziehung ist eine sehr gute Näherung, solange

\begin{equation}
	v
	\ll
	c
\end{equation}

gilt. Die Geschwindigkeit \(v\) des Teilchens muss also wesentlich
kleiner als die Lichtgeschwindigkeit \(c\) sein.

Für ein freies Teilchen lautet die Schrödingergleichung

\begin{equation}
	\mathrm{i}\hbar
	\frac{\partial}{\partial t}
	\Psi(\mathbf{r},t)
	=
	-\frac{\hbar^2}{2m}
	\nabla^2
	\Psi(\mathbf{r},t).
	\label{eq:freie_schroedingergleichung_grenzen}
\end{equation}

Der Ausdruck auf der rechten Seite entsteht unmittelbar aus der
klassischen Energieformel

\begin{equation}
	E_{\mathrm{kin}}
	=
	\frac{p^2}{2m}.
\end{equation}

Für hohe Geschwindigkeiten gilt jedoch die relativistische
Energie-Impuls-Beziehung%
\index{Energie-Impuls-Beziehung!relativistische}

\begin{equation}
	E^2
	=
	p^2c^2
	+
	m^2c^4.
	\label{eq:relativistische_energie_impuls_kapitel07}
\end{equation}

Diese Beziehung unterscheidet sich grundlegend von der klassischen
Formel. Insbesondere enthält sie die Ruheenergie%
\index{Ruheenergie}

\begin{equation}
	E_0
	=
	mc^2.
\end{equation}

\begin{PhysicsBox}[Nichtrelativistische Näherung]
	\label{box:nichtrelativistische_naeherung_schroedinger}
	
	Die Schrödingergleichung ist zuverlässig, solange die kinetische
	Energie des Elektrons wesentlich kleiner als seine Ruheenergie ist:
	
	\begin{equation}
		E_{\mathrm{kin}}
		\ll
		m_{\mathrm e}c^2.
	\end{equation}
	
	Für das Elektron beträgt die Ruheenergie
	
	\begin{equation}
		m_{\mathrm e}c^2
		\approx
		\SI{511}{\kilo\electronvolt}.
	\end{equation}
	
	Je näher die Energie des Elektrons an diese Größenordnung herankommt,
	desto wichtiger werden relativistische Korrekturen.
\end{PhysicsBox}

Das bedeutet nicht, dass die Schrödingergleichung bei jeder hohen
Energie plötzlich vollständig versagt. Relativistische Abweichungen
werden schrittweise größer. Für präzise Berechnungen müssen sie daher
bereits berücksichtigt werden, bevor die kinetische Energie die
Ruheenergie erreicht.

\subsubsection*{Raum und Zeit werden unterschiedlich behandelt}

In der Schrödingergleichung tritt die zeitliche Ableitung in erster
Ordnung auf:

\begin{equation}
	\frac{\partial\Psi}{\partial t}.
\end{equation}

Die räumliche Abhängigkeit erscheint dagegen über den
Laplace-Operator\index{Laplace-Operator} mit Ableitungen zweiter
Ordnung:

\begin{equation}
	\nabla^2\Psi
	=
	\frac{\partial^2\Psi}{\partial x^2}
	+
	\frac{\partial^2\Psi}{\partial y^2}
	+
	\frac{\partial^2\Psi}{\partial z^2}.
\end{equation}

Raum und Zeit werden mathematisch also unterschiedlich behandelt.

In der speziellen Relativitätstheorie bilden Raum und Zeit dagegen
eine gemeinsame vierdimensionale Raumzeit%
\index{Raumzeit}. Eine relativistische Gleichung muss so formuliert
sein, dass sie unter Lorentz-Transformationen%
\index{Lorentz-Transformation} ihre Form behält.

Die gewöhnliche Schrödingergleichung erfüllt diese Forderung nicht. Sie
ist an die nichtrelativistische Mechanik und deren
Galilei-Trans- \newline formationen%
\index{Galilei-Transformation} angepasst.

\begin{DidacticBox}[Was ist das eigentliche Problem?]
	\label{box:raum_zeit_schroedinger}
	
	Die Schrödingergleichung behandelt die Zeit anders als den Raum.
	
	Das ist im nichtrelativistischen Bereich kein Fehler. Dort ist diese
	Form genau die passende Näherung.
	
	Für eine relativistische Theorie reicht sie jedoch nicht aus, weil
	Raum und Zeit dann gemeinsam beschrieben werden müssen.
\end{DidacticBox}

\subsubsection*{Der Spin ist nicht Bestandteil der Gleichung}

Das Elektron besitzt einen intrinsischen Drehimpuls, den
Spin\index{Spin!Elektron}. Dieser Spin ist keine klassische Drehung
einer kleinen Kugel. Er ist eine grundlegende Quanteneigenschaft des
Elektrons.

Für das Elektron gilt

\begin{equation}
	s
	=
	\frac{\hbar}{2}.
	\label{eq:spin_elektron_grenzen}
\end{equation}

Die ursprüngliche Schrödingergleichung enthält den Spin nicht. Ein
Elektron wird darin zunächst nur durch eine skalare Wellenfunktion
beschrieben.

Der Spin kann nachträglich durch eine Erweiterung der Theorie
berücksichtigt werden. Dies führt zur Pauli-Gleichung%
\index{Pauli-Gleichung}. Dabei wird der Zustand des Elektrons durch zwei
Komponenten beschrieben, die den beiden möglichen Spinrichtungen
zugeordnet sind.

Die Pauli-Gleichung beschreibt auch die Wechselwirkung zwischen dem
Spin und einem Magnetfeld. Sie erklärt jedoch nicht, warum das Elektron
überhaupt den Spin \(\hbar/2\) besitzt.

\begin{DidacticBox}[Der Spin wird nicht hergeleitet]
	\label{box:spin_nicht_hergeleitet_schroedinger}
	
	Die Schrödingergleichung widerspricht dem Elektronenspin nicht.
	
	Der Spin ist jedoch nicht von Anfang an in ihrer mathematischen
	Struktur enthalten. Er muss zusätzlich eingeführt werden.
	
	Eine relativistische Theorie des Elektrons sollte den Spin dagegen
	unmittelbar aus der Gleichung selbst ergeben.
\end{DidacticBox}

\subsubsection*{Die Anzahl der Teilchen bleibt fest}

In der gewöhnlichen Quantenmechanik wird meist eine feste Anzahl von
Teilchen vorausgesetzt. Ein Elektron kann sich bewegen, gebunden werden
oder seinen Energiezustand ändern. Es bleibt dabei jedoch als Teilchen
erhalten.

Bei hohen Energien können Teilchen und Antiteilchen%
\index{Antiteilchen} jedoch erzeugt oder vernichtet werden.

Ein energiereiches Photon kann beispielsweise in der Nähe eines
Atomkerns ein Elektron-Positron-Paar erzeugen:

\begin{equation}
	\gamma
	\longrightarrow
	e^-
	+
	e^+.
	\label{eq:paarbildung_elektron_positron_grenzen}
\end{equation}

Dieser Vorgang wird als Paarbildung%
\index{Paarbildung!Elektron und Positron} bezeichnet.

Die dafür mindestens erforderliche Ruheenergie beträgt

\begin{equation}
	E_{\mathrm{min}}
	=
	2m_{\mathrm e}c^2
	\approx
	\SI{1.022}{\mega\electronvolt}.
	\label{eq:mindestenergie_paarbildung}
\end{equation}

Das Positron\index{Positron} besitzt dieselbe Masse wie das Elektron,
aber eine positive elektrische Ladung.

Eine Theorie mit einer fest vorgegebenen Teilchenzahl kann solche
Prozesse nicht vollständig beschreiben. Dafür wird später eine
Quantenfeldtheorie%
\index{Quantenfeldtheorie} benötigt.

\begin{PhysicsBox}[Feste Teilchenzahl]
	\label{box:feste_teilchenzahl_schroedinger}
	
	Die Schrödingergleichung beschreibt gewöhnlich Zustände mit einer
	festen Anzahl von Teilchen.
	
	Sie kann deshalb nicht vollständig erfassen, dass bei hohen Energien
	
	\begin{itemize}
		\item Teilchen erzeugt,
		\item Antiteilchen erzeugt und
		\item Teilchen wieder vernichtet werden können.
	\end{itemize}
	
	Solche Prozesse erfordern eine relativistische Quantenfeldtheorie.
\end{PhysicsBox}

\subsubsection*{Was die Schrödingergleichung dennoch leistet}

Die genannten Grenzen dürfen nicht den Eindruck erwecken, die
Schrödingergleichung sei nur eine grobe oder ungenaue Theorie.

Für nichtrelativistische Systeme ist sie außerordentlich erfolgreich.
Sie bildet unter anderem die Grundlage für die Beschreibung von

\begin{itemize}
	\item Elektronen in Atomen,
	\item chemischen Bindungen,
	\item Molekülen,
	\item Halbleitern,
	\item Festkörpern und
	\item zahlreichen quantenmechanischen Experimenten.
\end{itemize}

Ihre Grenzen treten vor allem dann hervor, wenn

\begin{itemize}
	\item sich Elektronen mit sehr hohen Geschwindigkeiten bewegen,
	\item starke relativistische Effekte auftreten,
	\item der Spin aus der Theorie selbst erklärt werden soll oder
	\item Teilchen erzeugt und vernichtet werden können.
\end{itemize}

\begin{DidacticBox}[Eine Näherung ist keine falsche Theorie]
	\label{box:naeherung_keine_falsche_theorie}
	
	Die Schrödingergleichung ist nicht falsch.
	
	Sie ist eine äußerst genaue Näherung für den nichtrelativistischen
	Bereich. Ihre Grenzen zeigen lediglich, dass sie nicht für alle
	Energien und alle physikalischen Prozesse geeignet ist.
	
	Auch die klassische Mechanik bleibt für langsame und große Körper
	richtig anwendbar, obwohl sie durch die Relativitätstheorie und die
	Quantenmechanik erweitert wurde.
\end{DidacticBox}

\subsubsection*{Die Suche nach einer neuen Gleichung}

Eine relativistische Gleichung für das Elektron musste mehrere
Bedingungen gleichzeitig erfüllen:

\begin{enumerate}
	\item Sie musste mit der speziellen Relativitätstheorie vereinbar
	sein.
	\item Sie musste die relativistische Energie-Impuls-Beziehung
	berücksichtigen.
	\item Sie musste eine sinnvolle Wahrscheinlichkeitsinterpretation
	ermöglichen.
	\item Sie sollte den Spin des Elektrons aus ihrer mathematischen
	Struktur erklären.
	\item Sie musste Lösungen zulassen, die mit Antiteilchen verbunden
	sind.
\end{enumerate}

Paul Dirac\index{Dirac, Paul} suchte nach einer Gleichung, die diese
Anforderungen erfüllt.

\begin{HistoryBox}[Der nächste Schritt]
	\label{box:naechster_schritt_dirac}
	
	Die Grenzen der Schrödingergleichung führten nicht zu ihrer Aufgabe,
	sondern zu einer Erweiterung der Quantentheorie.
	
	Paul Dirac verband die Quantenmechanik mit der speziellen
	Relativitätstheorie. Aus dieser Verbindung entstand eine neue Gleichung
	für das Elektron.
	
	Dabei ergaben sich der Spin und die Möglichkeit von Antiteilchen nicht
	als zusätzliche Annahmen, sondern als Folgen der mathematischen
	Struktur.
\end{HistoryBox}

Der Weg zu dieser Gleichung begann mit der Frage, wie sich die
relativistische Energie-Impuls-Beziehung in eine quantenmechanische
Bewegungsgleichung übertragen lässt.

